% LaTeX template for Artifact Evaluation V20201122
%
% Prepared by Grigori Fursin with contributions from Bruce Childers,
%   Michael Heroux, Michela Taufer and other colleagues.
%
% See examples of this Artifact Appendix in
%  * SC'17 paper: https://dl.acm.org/citation.cfm?id=3126948
%  * CGO'17 paper: https://www.cl.cam.ac.uk/~sa614/papers/Software-Prefetching-CGO2017.pdf
%  * ACM ReQuEST-ASPLOS'18 paper: https://dl.acm.org/citation.cfm?doid=3229762.3229763
%
% (C)opyright 2014-2023
%
% CC BY 4.0 license
%

\documentclass{sigplanconf}

\usepackage{hyperref}

\begin{document}

\special{papersize=8.5in,11in}

%%%%%%%%%%%%%%%%%%%%%%%%%%%%%%%%%%%%%%%%%%%%%%%%%%%%
% When adding this appendix to your paper, 
% please remove above part
%%%%%%%%%%%%%%%%%%%%%%%%%%%%%%%%%%%%%%%%%%%%%%%%%%%%

\appendix
\section{Artifact Appendix}

%%%%%%%%%%%%%%%%%%%%%%%%%%%%%%%%%%%%%%%%%%%%%%%%%%%%%%%%%%%%%%%%%%%%%
\subsection{Abstract}

{\em Obligatory}

\subsection{Artifact check-list (meta-information)}

{\em Obligatory. Use short keywords or one-line answers for the fields that
apply to your artifact, and remove the rest. This information is used to assign
appropriate reviewers. Detailed instructions can be placed in the artifact
README.}

{\small
\begin{itemize}
  \item {\bf Algorithm or method: }
  \item {\bf Programs, benchmarks, and data sets: }
  \item {\bf Models and model availability: }
  \item {\bf Compilation, binaries, tools, proprietary software, and PDKs: }
  \item {\bf Run-time environment, hardware, and run-time state: }
  \item {\bf Container, VM, or locked environment: }
  \item {\bf Workflow framework and execution: }
  \item {\bf Metrics, outputs, and tolerances: }
  \item {\bf Preparation/runtime, disk space, GPU/API cost: }
  \item {\bf Public availability, licenses, and restrictions: }
  \item {\bf Zenodo DOI for archived artifact: }
\end{itemize}
}

%%%%%%%%%%%%%%%%%%%%%%%%%%%%%%%%%%%%%%%%%%%%%%%%%%%%%%%%%%%%%%%%%%%%%
\subsection{Description}

\subsubsection{How to access}

{\em Obligatory}

\subsubsection{Hardware dependencies}

\subsubsection{Software dependencies}

\subsubsection{Commercial software and PDK dependencies}

\subsubsection{Data sets}

\subsubsection{Models}

For model-based artifacts, specify the exact model or API model, availability
or access instructions, version/revision where applicable, and license. If LLMs
materially affect the results, also summarize the prompts, inference settings,
metrics, raw or cached outputs, tolerance bands, and expected API cost or
local compute needed for validation.

%%%%%%%%%%%%%%%%%%%%%%%%%%%%%%%%%%%%%%%%%%%%%%%%%%%%%%%%%%%%%%%%%%%%%
\subsection{Installation}

{\em Obligatory}

%%%%%%%%%%%%%%%%%%%%%%%%%%%%%%%%%%%%%%%%%%%%%%%%%%%%%%%%%%%%%%%%%%%%%
\subsection{Experiment workflow}

%%%%%%%%%%%%%%%%%%%%%%%%%%%%%%%%%%%%%%%%%%%%%%%%%%%%%%%%%%%%%%%%%%%%%
\subsection{Evaluation and expected results}

{\em Obligatory}

For each key result, state the command or workflow reviewers should run, the
expected output or metric, the acceptable tolerance or similar-result criterion,
and whether the intended validation path is a full rerun, bounded subset rerun,
cached-output validation, or evidence/log inspection.

%%%%%%%%%%%%%%%%%%%%%%%%%%%%%%%%%%%%%%%%%%%%%%%%%%%%%%%%%%%%%%%%%%%%%
\subsection{Experiment customization}

%%%%%%%%%%%%%%%%%%%%%%%%%%%%%%%%%%%%%%%%%%%%%%%%%%%%%%%%%%%%%%%%%%%%%
\subsection{Notes}

%%%%%%%%%%%%%%%%%%%%%%%%%%%%%%%%%%%%%%%%%%%%%%%%%%%%%%%%%%%%%%%%%%%%%
\subsection{Methodology}

Submission, reviewing and badging methodology:

\begin{itemize}
  \item \url{https://www.acm.org/publications/policies/artifact-review-and-badging-current}
  \item \url{http://cTuning.org/ae/submission-20201122.html}
  \item \url{https://github.com/ml-eda/artifact-evaluation/}
\end{itemize}

%%%%%%%%%%%%%%%%%%%%%%%%%%%%%%%%%%%%%%%%%%%%%%%%%%%%
% When adding this appendix to your paper, 
% please remove below part
%%%%%%%%%%%%%%%%%%%%%%%%%%%%%%%%%%%%%%%%%%%%%%%%%%%%

\end{document}
